%
% board_ordering_current_buggy.tex
%
% FIDELITY CONTROL for docs/board_ordering.tex.  This is the applet's board
% path AS THE CODE STANDS -- the ungated push, the board captured before the
% raise, and the refresh that pushes before the slot has seen its board.  It
% exists to be violated: a model that cannot reproduce the defect it guards
% against is too abstract to trust.
%
% It differs from board_ordering.tex in exactly one section: the push.  The
% state, Init, TakePlace and StoreBoard are identical.
%
% Type-check:  fuzz -t docs/board_ordering_current_buggy.tex
% Model-check (each goal must be FOUND -- see Section 3):
%   probcli docs/board_ordering_current_buggy.tex -model_check \
%     -p DEFAULT_SETSIZE 2 -p MAXINT 4 -goal "shown /= evershown"
%

\documentclass[a4paper,10pt,fleqn]{article}
\usepackage[margin=1in]{geometry}
\usepackage{fuzz}
\usepackage[colorlinks=true,linkcolor=blue,citecolor=blue,urlcolor=blue]{hyperref}

\begin{document}

\title{Board Ordering as the Code Stands: a Fidelity Control}
\author{The variant that reaches the bad states}
\date{August 2026}
\maketitle

% ============================================================
\section{Purpose}
% ============================================================

This document is not a design. It is the control for
\texttt{docs/board\_ordering.tex}: the same state and the same slot, with the
push written as the code writes it today. Its job is to reach the bad states,
and it does.

The three staging operations below are the three pushes in the code.
\texttt{AnswerStage} is \texttt{HeldBoard.answer}, which is handed the board
\texttt{BeadsService.acknowledge} read out of the slot and pushes it after the
raise round trip has returned. \texttt{RefreshStage} is the push inside
\texttt{BoardRun.\_shown} and \texttt{NoBoard.refreshed}, which sends the board
this worker just loaded before \texttt{BoardSlot.store} has had the chance to
refuse it. \texttt{BlankStage} is \texttt{NoBoard.answer}'s placeholder and
\texttt{NoBoard.\_show\_reason}'s red message, which carry no board at all.

None of the three is gated: nothing in the code stops two pushes being in flight
at once, and nothing compares what a push is about to write against what the
glass is already showing.

% ============================================================
\section{Types, State, and the Slot}
% ============================================================

Identical to \texttt{docs/board\_ordering.tex}, which documents them.

\begin{zed}
  [WORKER]
\end{zed}

\begin{zed}
  PLACE == 0 \upto 3
\end{zed}

\begin{schema}{Board}
  began : PLACE \\
  loading : WORKER \pfun PLACE \\
  held : PLACE \\
  shown : PLACE \\
  staged : WORKER \pfun PLACE \\
  everheld : PLACE \\
  evershown : PLACE \\
  offered : PLACE \\
  overtaken : \power PLACE
\where
  held \leq began \\
  shown \leq began \\
  everheld \leq began \\
  evershown \leq began \\
  offered \leq began
\end{schema}

\begin{schema}{Init}
  Board'
\where
  began' = 0 \\
  loading' = \emptyset \\
  held' = 0 \\
  shown' = 0 \\
  staged' = \emptyset \\
  everheld' = 0 \\
  evershown' = 0 \\
  offered' = 0 \\
  overtaken' = \emptyset
\end{schema}

\begin{schema}{TakePlace}
  \Delta Board \\
  w? : WORKER
\where
  w? \notin \dom loading \\
  began < 3 \\
  began' = began + 1 \\
  loading' = loading \cup \{ w? \mapsto began' \} \\
  held' = held \\
  shown' = shown \\
  staged' = staged \\
  everheld' = everheld \\
  evershown' = evershown \\
  offered' = offered \\
  overtaken' = overtaken
\end{schema}

\begin{schema}{StoreBoard}
  \Delta Board \\
  w? : WORKER
\where
  w? \in \dom loading \\
  ( loading~w? > held \land held' = loading~w? \\
  \quad~ \land overtaken' = overtaken ) \lor \\
  \quad~ ( loading~w? \leq held \land held' = held \\
  \quad~ \land overtaken' = overtaken \cup \{ loading~w? \} ) \\
  ( held' > everheld \land everheld' = held' ) \lor \\
  \quad~ ( held' \leq everheld \land everheld' = everheld ) \\
  loading' = \{ w? \} \ndres loading \\
  began' = began \\
  shown' = shown \\
  staged' = staged \\
  evershown' = evershown \\
  offered' = offered
\end{schema}

% ============================================================
\section{The Push as the Code Writes It}
% ============================================================

\texttt{AnswerStage} is the click's answer: it takes the board the slot holds
and settles on pushing that one. In the code the read happens in
\texttt{acknowledge} and the push happens after \texttt{work.raise\_frame()}
has returned from a round trip to \texttt{luxd}, so the window between them is
wide. Here it is a step, which is all a window needs to be.

\begin{schema}{AnswerStage}
  \Delta Board \\
  w? : WORKER
\where
  w? \notin \dom staged \\
  staged' = staged \cup \{ w? \mapsto held \} \\
  ( held > offered \land offered' = held ) \lor \\
  \quad~ ( held \leq offered \land offered' = offered ) \\
  began' = began \\
  loading' = loading \\
  held' = held \\
  shown' = shown \\
  everheld' = everheld \\
  evershown' = evershown \\
  overtaken' = overtaken
\end{schema}

\texttt{RefreshStage} is the refresh's push, which sends the board this worker
loaded. It is enabled while that worker is still in \texttt{loading} --- that is
the defect: the push is settled before \texttt{StoreBoard} has compared this
board against the one the slot is keeping.

\begin{schema}{RefreshStage}
  \Delta Board \\
  w? : WORKER
\where
  w? \in \dom loading \\
  w? \notin \dom staged \\
  staged' = staged \cup \{ w? \mapsto loading~w? \} \\
  ( loading~w? > offered \land offered' = loading~w? ) \lor \\
  \quad~ ( loading~w? \leq offered \land offered' = offered ) \\
  began' = began \\
  loading' = loading \\
  held' = held \\
  shown' = shown \\
  everheld' = everheld \\
  evershown' = evershown \\
  overtaken' = overtaken
\end{schema}

\texttt{BlankStage} is a push carrying no board: the placeholder a click opens
when nothing has loaded, or the red message naming why \texttt{bd} could not be
read. It is enabled whenever the slot holds no board, and it consults the glass
about nothing.

\begin{schema}{BlankStage}
  \Delta Board \\
  w? : WORKER
\where
  w? \notin \dom staged \\
  held = 0 \\
  staged' = staged \cup \{ w? \mapsto 0 \} \\
  began' = began \\
  loading' = loading \\
  held' = held \\
  shown' = shown \\
  everheld' = everheld \\
  evershown' = evershown \\
  offered' = offered \\
  overtaken' = overtaken
\end{schema}

\texttt{PushCommit} is the write landing, exactly as in the committed model.

\begin{schema}{PushCommit}
  \Delta Board \\
  w? : WORKER
\where
  w? \in \dom staged \\
  shown' = staged~w? \\
  ( staged~w? > evershown \land evershown' = staged~w? ) \lor \\
  \quad~ ( staged~w? \leq evershown \land evershown' = evershown ) \\
  staged' = \{ w? \} \ndres staged \\
  began' = began \\
  loading' = loading \\
  held' = held \\
  everheld' = everheld \\
  offered' = offered \\
  overtaken' = overtaken
\end{schema}

% ============================================================
\section{What it reaches}
% ============================================================

The same five goals as the committed model. Here four of them are
\emph{found}, which is the point of the document.

\begin{verbatim}
  probcli docs/board_ordering_current_buggy.tex -model_check \
    -p DEFAULT_SETSIZE 2 -p MAXINT 4 -p TIME_OUT 60000 \
    -goal "held /= everheld"                 # I1: NOT found -- the slot is sound
  probcli ... -goal "shown /= evershown"              # I2: FOUND
  probcli ... -goal "staged = {} & shown < offered"   # I3: FOUND
  probcli ... -goal "shown > held"                    # I4: FOUND
  probcli ... -goal "shown : overtaken"               # I5: FOUND
\end{verbatim}

I1 not found is a result in its own right: the slot's own ordering rule was
never the problem, and none of these defects is a defect in
\texttt{BoardSlot}. The other four are the defect, in four shapes:

\begin{description}
  \item[I2, the glass goes backwards.] A click's answer settles on the board the
    slot holds; a refresh stores and lands a newer one; the click lands its
    older capture over it. If the click then stood down --- and it will have, if
    the refresh that overtook it is the one holding single flight --- nothing
    repairs the display.
  \item[I3, the newest board offered never reaches the glass.] The same
    interleaving seen from the other side: a board at place $2$ was settled on
    and the display came to rest at place $1$.
  \item[I4, the glass runs ahead of the slot.] \texttt{RefreshStage} followed by
    \texttt{PushCommit} with no \texttt{StoreBoard} between: the display is
    rendering a board the slot has not accepted and may yet refuse.
  \item[I5, the glass shows a refused board.] Two loads begin, the later one
    stores first, and the earlier one lands its board on the glass and is then
    refused by the slot. The applet holds one board and displays another, and
    the one it displays is the one it decided against.
\end{description}

Over 3275 reachable states, I1 is not found and the other four are found in
seven to eleven steps.

% ============================================================
\section{Realisability}
% ============================================================

This model omits \texttt{SingleFlight}, so it admits interleavings the running
applet cannot produce --- two refreshes at once, in particular. That direction
is safe for the committed design, where the claim is that no trace violates
anything. It is not safe here, where the claim is that a particular trace
exists: a witness that needs two concurrent refreshes would be an artifact.
So each is checked against the code by hand.

\paragraph{I2, a placeholder over a board (seven steps).}
Click A takes the single-flight token and starts a load. Click B arrives on
another thread, reads the slot in \texttt{acknowledge}, finds \texttt{NoBoard},
and settles on the placeholder; it then spends the raise round trip. A's load
returns, pushes its board, and stores it. B's placeholder lands. B finds the
token held and stands down, so nothing runs behind it to repair the display.
The user clicked twice and is reading ``Loading issues\ldots'' over a board that
loaded. One refresh, one answer --- realisable.

\paragraph{I2, a board over an older board (nine steps).}
The same shape with a warm slot. Click B's \texttt{acknowledge} captures the
board at place $1$ and goes into its raise. Click C, on another thread, holds
the token, loads place $2$, pushes it, and stores it. B's captured $1$ then
lands. B stands down. The applet holds place $2$ and the display shows place
$1$, permanently. One refresh, one answer --- realisable. This is the defect as
reported.

\paragraph{I4, the glass ahead of the slot (five steps).}
No race at all is needed. \texttt{BoardRun.\_shown} pushes the board and
\texttt{BeadsService.service} stores it afterwards, so between those two
statements the display is showing a board the slot has not accepted. Every
refresh passes through that state; the model finds it in one worker.

\paragraph{I5, a refused board on the glass (six steps).}
Needs a store the pusher does not control, which the warm-up provides: it
stores and never pushes. A click starts its load; a reconnect fires the warm-up
on another thread; the warm-up's load begins later, so it holds the newer
issues, and it stores first. The click's load returns, pushes its own older
board, and is then refused by the slot. The applet holds the newer board, the
display shows the one it refused, and no further click is pending to correct it.
One refresh and one warm-up --- realisable, and the warm-up-after-reconnect path
is the one the warm-up exists for.

\end{document}
